\documentclass{unbthesis}
\usepackage{amssymb,amsfonts,amstext,amsmath}
\DeclareMathOperator{\sen}{sen}
\DeclareMathOperator{\argum}{arg}
\usepackage[brazil]{babel}
\usepackage{blindtext}
\usepackage[T1]{fontenc}
\usepackage[utf8]{inputenc} 
\usepackage{color}
\usepackage{subfigure}
\usepackage{setspace}
\usepackage{longtable}
\usepackage{float}
\usepackage{colortbl}
\usepackage{verbatim}
\usepackage{lscape}
\usepackage{indentfirst}
\usepackage[bookmarks=true,bookmarksopen=false,linktocpage=true,pagebackref]{hyperref}
\usepackage[alf,abnt-etal-text=it,abnt-etal-list=0,abnt-etal-cite=2,abnt-and-type=&,abnt-year-extra-label=yes]{abntex2cite}

\usepackage{algorithmic}
\usepackage{textcomp}
\usepackage{stfloats}
\usepackage{xfrac}
\usepackage{booktabs}
\usepackage{soul}
\usepackage{subfigure}
\usepackage{multirow}
\usepackage{hhline}
\usepackage{graphics,graphicx}
\usepackage{epstopdf}
\usepackage{array}
\usepackage{icomma}

% Configurações do pacote backref
% Usado sem a opção hyperpageref de backref
\renewcommand{\backrefpagesname}{Citado na(s) página(s):~}
% Texto padrão antes do número das páginas
\renewcommand{\backref}{}
% Define os textos da citação
\renewcommand*{\backrefalt}[4]{
\ifcase #1 %
Sem citações no texto.%
\or
Citado na página #2.%
\else
Citado #1 vezes nas páginas #2.%
\fi}%


\usepackage{booktabs}
\usepackage{ctable}
\usepackage{multirow}
\usepackage{fancyhdr}
\usepackage{fancyvrb}
\usepackage{relsize}
\usepackage{calc}
\usepackage{psfrag}
\usepackage{epstopdf}
\usepackage{wallpaper}
\usepackage{paralist}
%\usepackage{hhline} %pacote para criar linhas horizontais segmentadas
%\usepackage{url}
%\usepackage{layout}
%\usepackage[retainorgcmds]{IEEEtrantools}
\usepackage{afterpage}

\setlength{\LTcapwidth}{17cm} \setlength{\textwidth}{17cm}
\setlength{\textheight}{24.9cm} \setlength{\topmargin}{-0.5in}

%NOVOS COMANDOS
\newcommand{\FuncSen}[1]{\mbox{\hspace{0.5mm}sen\hspace{0.5mm}}#1}
\newcommand{\FuncCos}[1]{\mbox{\hspace{0.5mm}cos\hspace{0.5mm}}#1}
\newcommand{\palfa}[1]{\mbox{plano-$\alpha$}#1}

\usepackage{mathtools}
\DeclarePairedDelimiter\abs{\lvert}{\rvert}%

%DEIXA ESPAÇOS EXTRAS, QUE O LATEX CRIA, NO FIM DA PÁGINA
%\raggedbottom



%\addto\captionsbrazil{%
%	\renewcommand{\contentsname}{Table of Contents}
%	\renewcommand{\listfigurename}{List of Figures}
%	\renewcommand{\listtablename}{List of Tables}
%	\renewcommand{\refname}{References}
%	\renewcommand{\bibname}{References}
%	\renewcommand{\chaptername}{Chapter}
%	\renewcommand{\figurename}{Figure}
%	\renewcommand{\tablename}{Table}
%	\renewcommand{\acknowledgementsname}{Acknowledgements}
%}


%%%VARIÁVEIS DO TEMPLATE - AJUSTAR DE ACORDO - se usar direito vai ter menos trabalho!

%variáveis exibidas na capa e contracapa
\newcommand{\ttitle}{Template \LaTeX para dissertações e teses do ENE/PPGEE/UnB}
\newcommand{\tauthor}{Comissão de internacionalização ENE/PPGEE/UnB}
\newcommand{\tdoctype}{Trabalho de Conclusão de Curso}
% ou \newcommand{\tdoctype}{Tese de Doutorado}
\newcommand{\tcourse}{Engenharia Elétrica}
\newcommand{\tdepartment}{Departamento de Engenharia Elétrica}
\newcommand{\tdpt}{ENE} %ou ENC ENM etc

\newcommand{\tdegree}{Engenheiro(a) Eletricista}
% ou \newcommand{\tdegree}{Doutor}

%variáveis utilizadas na ficha catalográfica
\newcommand{\tlocal}{Brasília/DF}
\newcommand{\tday}{03}
\newcommand{\tmonth}{setembro}
\newcommand{\tyear}{2022}

%nome para citação
\newcommand{\tsurname}{Sobrenome}
\newcommand{\tfirstname}{Nome}



\begin{document}

%PEDIDO DE MODELO COM AS DIMENSÕES DA PÁGINA
%PRECISA DO PACKAGE "layout"
%\layout

%NOVO ESPAÇAMENTO
\baselineskip 20pt
\parskip 5pt
\pagestyle{empty}

\pdfbookmark[0]{Capa}{Capa}
\include{01_cover}
\include{02_approval}
\include{03_copyright}

%NOVO ESPAÇAMENTO
\baselineskip 22pt
\pdfbookmark[0]{Dedicatória}{Dedicatória}
\include{03.1_dedicatory}
\pdfbookmark[0]{Agradecimentos}{Agradecimentos}
\include{03.2_acknowledgements}
\pdfbookmark[0]{Resumo}{Resumo}
\include{03.4_resumo}
\pdfbookmark[0]{Abstract}{Abstract}
\include{03.3_abstract}
%% DEFINE CABEÇALHO DO SUMÁRIO E DAS LISTAS

%IMPRESSÃO SOMENTE FRENTE
\pagestyle{fancy}
\renewcommand{\chaptermark}[1]{\markboth{ #1}{}}
\fancyhf{}
\fancyhead[RE,RO]{\fontsize{9}{10}\thepage}
\fancyhead[LO,LE]{\fontsize{9}{10}\textit{\textsc{\nouppercase\leftmark}}}

%NOVO ESPAÇAMENTO
\baselineskip 14pt
\pagenumbering{roman}

%% Sumário (gerado automáticamente)
\phantomsection %comando para fazer os bookmarks funcionarem na seção abaixo
\addcontentsline{toc}{chapter}{Sumário}
\tableofcontents

%NOVO ESPAÇAMENTO
\baselineskip 23pt

%% Lista de figuras (gerada automaticamente)
\clearpage
\phantomsection
\addcontentsline{toc}{chapter}{Lista de figuras}
\listoffigures

%% Lista de tabelas (gerada automaticamente)
\clearpage
\phantomsection
\addcontentsline{toc}{chapter}{Lista de tabelas}
\listoftables

%% Lista de Simbolos (gerada manualmente)
\markboth{LISTA DE SÍMBOLOS}{LISTA DE SÍMBOLOS}
\clearpage
\phantomsection
\addcontentsline{toc}{chapter}{Lista de símbolos}
\include{10_symbols}

%% Glossario (gerada manualmente)
\markboth{GLOSSÁRIO}{GLOSSÁRIO}
\clearpage
\phantomsection
\addcontentsline{toc}{chapter}{Glossário}
\include{11_glossary}

\pagestyle{empty}
\pagenumbering{arabic}

%NOVO ESPAÇAMENTO
\baselineskip 23pt
\parskip 5pt

%IMPRESSÃO SOMENTE FRENTE
\pagestyle{fancy}
\renewcommand{\chaptermark}[1]{\markboth{\thechapter\hspace{2mm}--\hspace{2mm}#1}{}}
\renewcommand{\sectionmark}[1]{\markright{\thesection\hspace{2mm}--\hspace{2mm}#1}{}}
\fancyhf{}
\fancyhead[RE,RO]{\fontsize{9}{10}\textit{\textsc{\thepage}}}
\fancyhead[LE,LO]{\fontsize{9}{10}\textit{\textsc{\rightmark}}}
\renewcommand{\footrulewidth}{0.0pt} %Definindo a espessura de uma linha do rodapé

\include{04_cp1_introduction}
%\include{05_cp2_foundations}
%\include{06_cp3_bibreview}
%\include{07_cp4_circuitapproach}
%\include{08_cp5_emapproach}
%\include{09_cp6_conclusions}

%\include{11-Cap_3_FundamentacaoTeorica}
%\include{10-Cap_2_RevisaoBibliografica}
%\include{13-Cap_5a_GM}
%\include{12-Cap_4_MelhoriasPropostas}
%\include{13-Cap_5_ApresAnalDosResultados}
%\include{14-Cap_6_ConclusaoPropostas}

% DEFINE CABEÇALHO PARA A BIBLIOGRAFIA
\pagestyle{fancy}
\fancyhead[RE,RO]{\fontsize{9}{10}\textit{\textsc{\thepage}}}
\fancyhead[LE,LO]{\fontsize{9}{10}\textit{\textsc{\nouppercase{\rightmark}}}}

%% BIBLIOGRAFIA (GERADO AUTOMATICAMENTE)
\clearpage
\phantomsection
\addcontentsline{toc}{chapter}{Referências}
\bibliography{library}

% DEFINE CABEÇALHO PARA O APÊNDICE
\pagestyle{fancy}
\fancyhead[RE,RO]{\fontsize{9}{10}\textit{\textsc{\thepage}}}
\fancyhead[LE,LO]{\fontsize{9}{10}\textit{\textsc{\nouppercase{\rightmark}}}}

%\appendix
%\include{15-Apendices}

\end{document} 